\documentclass{article}
\usepackage[utf8]{inputenc}
\usepackage{ulem}
\usepackage{cancel}
\title{Math Quiz 1}
\author{Cole Fierst }
\date{September 12th 2021}

\begin{document}

\maketitle

\section{Problem 1}


D is defined in terms of y, so the equation for this parabola has to be:

$$ y-c=r(x-b)^2 $$
We are given the focus and directrix
$$ F : (x,y) = (2,-1),\hspace{3}D : y=3$$
And the equation's focus and directrix shifts are
$$  F=\left(b,c+\frac{1}{4r}\right),\hspace{3} D:y=c-\frac{1}{4r}$$
Plug in relevant values from given numbers into equation's F and D
$$ 2=b,\hspace{4} -1=c+\frac{1}{4r}, \hspace{4} 3=c-\frac{1}{4r}$$
Make values equal by adding 4 to equation that was equal to -1
$$ 3=c+\frac{1}{4r}+4$$
System of equations to solve for r
$$ c+\frac{1}{4r}+4=c-\frac{1}{4r} $$
$$\frac{2}{4r} =-4$$

$$ \frac{1}{2r}=-4 $$

$$ 1 = -8r $$
$$ r=-\frac{1}{8} $$
Plug in r to solve for c
$$ D:3=c-\frac{1}{4(-\frac{1}{8})}\hspace{3}\rightarrow\hspace{3}3=c--\frac{1}{1/2}$$
$$ 3=c+2$$
$$ c=1$$
Have everything to complete equation giving us a final answer of:
$$ y-1 = -\frac{1}{8}(x-2)^2 $$


\section{Problem 2}

The inequality can also be expressed as
$$ f_{\alpha}(4)>f_{10}(10)$$
$\alpha>0$ for all observed instances of $\mathrm{4^\alpha}$ 
\\
This means that since $f_{\alpha}(10)$ where $\alpha=10$ that both functions are \textbf{Strictly increasing} since $\alpha > 0$ is always true
\\
\\
\textbf{a)}
\sout{Not large enough, if we use the laws of exponents, we can multiply both sides by 10th root}
$$ \cancel{\displaystyle4^{10*\frac{1}{10}}>10^{10*\frac{1}{10}}} $$
\sout{Which gives us a false inequality}
$$ \cancel{ 4>10} $$
\\
\textbf{REVISED A}
\\
If $n=10$, then inequality is 
$$4^{10}>10^{10}$$
Since both sides are increasing, and to the same power, we can compare bases directly
$$4>10$$
Which is an false inequality, meaning $n=10$ isn't enough
\\
\\
\textbf{b)}
Yes it's large enough, we can seperate $4^{20}$ by using law of exponents
$$ 4^{20}=(4^{10})^2$$
$$ (4^{10})^2=(4^2)^{10} $$
$$(4^2)^{10}=16^{10}$$
\sout{Apply same principle as in a, but being careful to distribute $\frac{1}{10}$ power equally}
$$ \cancel{4*4>10} $$
Now that we have both sides of inequality to same power, we can again compare bases since both are also increasing
$$ 16^{10}>10^{10} \rightarrow 16>10 $$
Inequality is true, so n=20 is enough
\\
\\
\textbf{c)}
\sout{No it isn't. By breaking up both sides of equation into 5th degree exponents, we can prove $4^{15}$ isn't big enough while making sure the bases and exponents stay as \textit{whole numbers}}
$$ \cancel{ 4^5*4^5*4^5=10^5*10^5} $$
\sout{Multiply both sides by $\frac{1}{5}$ power, making sure to distribute correctly}
$$ \cancel{4^3>10^2} $$
\\
$n=15$ isn't enough. By applying same principal as other two, we are going to use laws of exponents to get get both sides to a common exponent
$$4^{15}>10^{10}$$
This time have to rearrange both sides, this time using $5$ since it's a common multiple between $15$ and $10$
$$(4^5)^3>(10^5)^2$$
$$(4^3)^5>(10^2)^5$$
$$64^5>100^5$$
Now if we compare bases after getting to common exponent when both are strictly increasing, we get:
$$ 64>100 $$
This inequality is false, so $n=15$ isn't enough

\section{Problem 3}
\textbf{a)}
$4x^3$, previous examples follow a $nx^{n-1}$ pattern, but we haven't proved this true at time of the quiz (but my future self correcting this has now seen it proven)
\\
\textbf{b)}
Any value $s$ in relation to a fixed value $x$ can be compared vertically by
$$ s^4-x^4 $$
Unlike above (when compared to ones for previous degrees), the tangent line equation doesn't change
$$ m(s-x) $$
We want the line to always be below the parabola, so for any value s, we want
$$ s^4 - x^4 \geq m(s-x) $$
Move over $m(s-x)$ and factor out common $(s-x)$ gives us
$$ (s^3+s^2x+sx^2+x^3-m)(s-x)\geq 0 $$
Have to "know" for geometric argument, fortunately, we do know
$$ m=4x^3 $$
$$ (s^3+s^2x+sx^2-3x^3)(s-x)\geq 0 $$
\sout{Parts of expression cancel out if multiplied by$(s-x)$}
\\
$$\cancel{{(s^4-4x^3s+3x^4)\geq 0}}$$
\\
Instead of above, will realize that $s^3+s^2x+sx^2-3x^3$ CAN be factored by $(s-x)$.
$$ (s-x)(s^3+s^2x+sx^2-3x^3)=(s-x)\mathbf{(s-x)(s^2+2sx+3x^2)}$$
Fact check that $(s-x)$ is a factor
$$ (s-x)(s^2+2sx+3x^2)= (s^3+2s^2x+3x^2s) + (-s^2x-2x^2s-3x^3)$$
$$ s^3+(2s^2x-s^2x)+(3x^2s-2x^2s)-3x^3=s^3+s^2x+x^2s-3x^3$$
Math above checks out, so we carry on
\\
Can factor even further since $s^2+2sx+3x^2$ can be factored again
$$ s^2+2sx-3x^2=(s+x)(s+x)+2x^2$$
After all this factoring we are left with
$$ (s-x)(s-x)(s+x)(s+x)+2x^2=(s-x)^2(s+x)^2+2x^2 $$
Every factor is to 2nd power, so are \textbf{all} are non-negative, meaning:
$$(s-x)^2(s+x)^2+2x^2 \geq 0$$
Which proves $4x^3=m$

\sout{I couldn't figure out where to go next in my 4 thinking hours. $s^4$ and $3x^4$ are always non negative, but $-4x^3s$ can be anything. Because of that, I wasn't sure how I would prove that the inequality is always greater than or equal to 0.  If the shaky limits argument was applied to this as it was, it would be $x^4-4x^4+3x^4$ which is just zero}


\section{Problem 4}
For fixed points $P (a,a^2)$ and $Q(b,b^2)$, we can connect them to an unfixed point $R(c,c^2)$ between them somewhere on the parabola
\\
$\tilde{R}$ is a possible Point $R$ where the triangle from the 3 points will have the greatest area of any in the parabola
\\
The point $\tilde{R}$ is located exactly between $P$ and $Q$ represented by $\tilde{R} = \frac{a+b}{2}$
\\
$\tilde{R}$ can equal $R$, but not usually
\\
The slope of the tangent line of $\tilde{R} = a+b$ and it is parallel to Line $PQ$
\\
For all triangles:
$$ A_{tri}PQR \leq A_{tri}P\tilde{R}Q$$
We can replace $A_{tri}$ with actual formula for area of a triangle to compare which variables control it
\\
Let $\tilde{H}$ represent height of the triangle with $\tilde{R}$
\\
$$ \frac{PQ*H}{2}\leq\frac{PQ*\tilde{H}}{2} $$
after canceling out constants:
$$ H \leq \tilde{H}$$
the only variable that matter for the area of a triangle in the parabola is height
\\
$\tilde{R}$ is also the only point on the parabola that is also on tangent $\stackrel{\leftrightarrow}{L}$
\\
This makes sense given that tangent lines also intersect what they are tangent to \textit{once} (in the general area, which in case of 2nd degree parabola is everwhere since it won't meet the parabola \textit{ever} again) 
\begin{center}
Otherwise, the line is below the parabola
\end{center}
\\
\begin{center}
$\stackrel{\leftrightarrow}{L}$  $\parallel$ $\stackrel{\leftrightarrow}{PQ}$ since they have the same slope
\end{center}

\\
Using these premises, and knowing that height is the only indicator of $A_{tri}$, we can say about height:
\begin{center}
H = Distance from $\stackrel{\leftrightarrow}{PQ}$ to point $R$
\\
Additionally, the distance between $\stackrel{\leftrightarrow}{PQ}$ and $\stackrel{\leftrightarrow}{L}$ = Distance from $\stackrel{\leftrightarrow}{PQ}$ to point $R$ (which equals $H$) $+$ distance from point $R$ to $\stackrel{\leftrightarrow}{L}$
\\
However, $\tilde{R}$ is on $\stackrel{\leftrightarrow}{L}$, so the distance from point $\tilde{R}$ to $\stackrel{\leftrightarrow}{L}$ is 0
\\
This means the distance between $\stackrel{\leftrightarrow}{PQ}$ and $\stackrel{\leftrightarrow}{L}$ = $H + 0$, meaning the distance between the parallel lines equals $H$ for, and only for $\tilde{R}$
\\
This is also the greatest height since the maximum height is the distance between $\stackrel{\leftrightarrow}{PQ}$ and $\stackrel{\leftrightarrow}{L}$, which $H$ is when solved for $\tilde{R}$
\\
Since height is the only relevant variable to the area of the triangle, the greatest height at ${\tilde{R}}$ is also the greatest area, proving why Archimedes chose ${\tilde{R}}$ for the 3rd point 
\end{center}
\sout{I will end my write up on 4 here since I ran out of time and was likely not going to progress further in the problem if I spent more time. Had thought in head of using the parallel lines of $PQ$ and tangent line $\tilde{R}$ to create a rectangle, but ultimately didn't see how that would allow me to find relationship between area.}

\end{document}
